\documentclass[../DoAn.tex]{subfiles}
\begin{document}

Chương này giới thiệu bài toán lập kế hoạch bao phủ cho robot tự hành trong môi trường lưới có ràng buộc năng lượng và vật cản động. Bối cảnh ứng dụng được chọn là robot mặt đất hoạt động trong khu vực nguy hiểm, chẳng hạn rà soát vật thể nghi ngờ hoặc khảo sát khu vực khó tiếp cận. Báo cáo không mô phỏng cảm biến dò mìn hay phần cứng robot thật; trọng tâm là mô hình hóa, lập kế hoạch đường đi, quản lý năng lượng và đánh giá hành vi bằng dữ liệu mô phỏng.

\section{Đặt vấn đề}

Trong bài toán bao phủ, robot không chỉ đi từ một điểm đầu tới một điểm đích. Nó phải lần lượt tiếp cận nhiều ô trong vùng làm việc, đánh dấu các ô đã được kiểm tra và duy trì trạng thái nhiệm vụ an toàn. Với robot mặt đất trong môi trường nguy hiểm, một đường đi ngắn về mặt hình học chưa chắc là lựa chọn tốt nếu đường đó tiêu tốn nhiều năng lượng, buộc robot đổi hướng liên tục hoặc làm robot khó quay về căn cứ.

\begin{figure}[H]
\centering
\begin{subfigure}{0.48\textwidth}
    \centering
    \includegraphics[width=\textwidth]{Hinh_ve/robot_talon.png}
    \caption{Robot TALON}
    \label{fig:robot_talon}
\end{subfigure}
\hfill
\begin{subfigure}{0.48\textwidth}
    \centering
    \includegraphics[width=\textwidth]{Hinh_ve/robot_packbot_510.png}
    \caption{Robot PackBot 510}
    \label{fig:robot_packbot_510}
\end{subfigure}
\caption{Ví dụ các robot mặt đất cỡ nhỏ hoạt động trong môi trường nguy hiểm. Nguồn ảnh: Army Technology.}
\label{fig:mine_robot_illustration}
\end{figure}

Hình~\ref{fig:mine_robot_illustration} chỉ dùng để minh họa động cơ ứng dụng. Trong phạm vi đồ án, robot được trừu tượng hóa thành một tác nhân di chuyển trên bản đồ lưới. Mỗi ô có thể là ô hợp lệ, ô vật cản hoặc ô có chi phí địa hình riêng. Căn cứ là vị trí xuất phát và cũng là nơi robot quay về để sạc hoặc kết thúc nhiệm vụ.

Ba ràng buộc làm bài toán trở nên khác với tìm đường ngắn nhất thông thường. Thứ nhất, robot có năng lượng hữu hạn nên không thể chỉ chọn mục tiêu gần nhất mà bỏ qua đường quay về. Thứ hai, hướng hiện tại của robot ảnh hưởng đến chi phí, vì thao tác đổi hướng có thể tiêu tốn năng lượng. Thứ ba, môi trường có thể thay đổi theo thời gian khi có vật cản động, buộc robot phải kiểm tra an toàn và lập kế hoạch lại khi cần.

\section{Mục tiêu của đồ án}

Mục tiêu của đồ án là xây dựng một hệ thống mô phỏng robot bao phủ bản đồ lưới có trọng số. Hệ thống cần đáp ứng các yêu cầu sau:

\begin{enumerate}
\item Biểu diễn môi trường bằng bản đồ lưới có vật cản tĩnh và chi phí địa hình.
\item Lập kế hoạch đường đi có xét vị trí và hướng hiện tại của robot.
\item Tính năng lượng tiêu thụ gồm năng lượng di chuyển và năng lượng quay.
\item Chỉ chọn mục tiêu bao phủ khi robot còn đủ năng lượng để quay về căn cứ.
\item Xử lý vật cản động ở mức mô phỏng thông qua kiểm tra ô an toàn và lập kế hoạch lại.
\item Ghi log, xuất benchmark và trực quan hóa kết quả để phục vụ đánh giá thực nghiệm.
\end{enumerate}

Các yêu cầu trên được giữ ở mức mô phỏng thuật toán. Những vấn đề như cảm biến thực, bản đồ xác suất, sai số định vị, điều khiển liên tục và thiết kế cơ khí nằm ngoài phạm vi chính của báo cáo.

\section{Đóng góp của đồ án}

Điểm chính của đồ án không nằm ở việc đề xuất một thuật toán tìm đường hoàn toàn mới. Đóng góp nằm ở cách tích hợp các thành phần đã biết thành một mô hình mô phỏng có thể chạy, quan sát và đánh giá định lượng.

Cụ thể, mô hình đề xuất khác với mô phỏng bao phủ chỉ dựa trên số bước ở ba điểm. Một là trạng thái tìm đường được mở rộng từ vị trí sang cặp vị trí--hướng, nhờ đó chi phí quay được đưa vào quá trình lập kế hoạch. Hai là điều kiện quay về căn cứ được kiểm tra trước khi robot chọn mục tiêu mới, giúp nhiệm vụ không chỉ tối đa hóa coverage mà còn xét khả năng thu hồi robot. Ba là hệ thống có cơ chế ghi benchmark theo nhiều nhóm kịch bản, cho phép so sánh ảnh hưởng của vị trí căn cứ, năng lượng tối đa, mật độ vật cản, cấu trúc vật cản, địa hình có trọng số và mô hình toàn hướng lý tưởng.

\section{Phạm vi và giả thiết}

Nghiên cứu này dùng một robot duy nhất trên bản đồ lưới hai chiều đã biết. Robot di chuyển theo bốn hướng chính và mỗi bước chỉ đi sang một ô kề cạnh. Vật cản động được xem là tác nhân không đối kháng; chúng làm thay đổi trạng thái an toàn của một số ô nhưng không chủ động truy đuổi robot. Mô hình năng lượng là mô hình chuẩn hóa, dùng để so sánh đường đi trong mô phỏng thay vì đo trực tiếp năng lượng vật lý của một robot cụ thể.

Phạm vi này giúp báo cáo tập trung vào lõi của bài toán Coverage Path Planning có ràng buộc năng lượng. Các hướng mở rộng như nhiều robot, cảm biến dò tìm, bản đồ xác suất hoặc điều khiển liên tục được xem là hướng phát triển sau.

\section{Bố cục đồ án}

Chương 2 trình bày cơ sở lý thuyết: bản đồ lưới, đồ thị có trọng số, trạng thái có hướng, chi phí quay, Dijkstra có hướng, mô hình năng lượng và tập ô an toàn. Chương 3 mô tả phương pháp đề xuất và thiết kế hệ thống. Chương 4 phân tích các lựa chọn mô hình và độ phức tạp. Chương 5 đánh giá thực nghiệm bằng benchmark. Chương 6 tổng kết kết quả và nêu hướng phát triển.

\end{document}
